\subsection{封国不是一排听令的据点}

\begin{event}{周初至成康时期}{重组东方封国网络}{总体过程可靠；各国始封的年月、受封者与边界常有争议}{成周以东及各地政治节点}\eventanchor{WZ-EAST-ENFEOFFMENT}{西周·重组东方封国网络}
东征的军队撤去之后，东方留下的不是一张可以交接的地图，而是许多仍在运作的宗族、旧商贵族、聚落与交通通道。周王室必须决定：哪些人可被托付，哪些旧有关系应被拆开，哪些地方要保留足以集结粮食、兵员和祭祀的政治节点。宗亲、功臣和被重新纳入的旧商贵族因此陆续进入不同位置；后世所说的鲁、齐、燕、卫、宋等封国，正是在这段漫长而不齐整的重组中获得了可继承的政治名分。

传世叙事喜欢把这些安排写成一连串清楚的授封：谁在某年受土，谁在何处立国，疆界又到哪里。然而，能较早看见的材料并不提供这样的总册。\sourcebook{尚书·康诰}、\sourcebook{酒诰}与\sourcebook{梓材}保存了对康叔及殷民的告诫和治理语言，却不是附有地籍的封授文书；\sourcebook{史记}各世家把鲁、齐、燕、卫、宋的祖源串成较完整的谱系，成书却远在事件之后。鲁侯簋、燕侯相关金文和琉璃河等地的早周遗存，能让贵族礼器、墓地和地方经营从纸面走到地下，却仍不能替每一国签出确切的始封日期、第一道边界或一份移民名册。

因此，较稳妥的说法不是“周初一次分完了东方”，而是王室在\eventref{WZ-1040-DONGZHENG}{西周·三监东征}之后，以册命、婚姻、驻守、祭祀和资源调动反复确认若干地方节点。对受命家族而言，这带来了可传给后代的地位；对当地原有共同体而言，则意味着须在新来的政治中心、旧有生计与服役要求之间重新安排关系。它的直接后果，是东方不再只靠一支远征军维持联系；它的中期代价，是王室必须不断用赏赐、裁断和亲临，使这些各有根基的节点仍愿接续王命。
\end{event}

成周的建设必须和东方封国的布置一起看。后来常用“分封制”概括周初安排；这一概念有助于说明王室把土地、名分和防务交给不同政治集团，却容易让人误以为当时已经有一张整齐的行政区划图。周王直接掌握的王畿、宗室支系的封地、功臣和盟友的政治基础，以及原有地方共同体之间，实际存在不同层次的关系。各地形成的先后、疆界与王室控制程度也并不相同。

王室提供的不是单纯的一块土地，而是一套可被承认的资格：受命者可以在当地组织宗族、祭祀和武装，并以周王授予的名分来处理对内对外关系。作为交换，王室期待他们在朝见、用兵、礼仪和资源供给上承担义务。这里的“期待”尤其重要。西周没有一套留存至今、可逐条核对的统一地方税法；不同封国在何时以何种物资、兵员或礼仪回应王室，不能用后世官僚制的标准倒推。

这是一种低成本却不廉价的统治方法。王室不必在每一处都长期派驻同等规模的官员和军队，封国的宗族资源能够承担守卫、开垦、调解与扩张；但中心要持续确认其地位，处理彼此冲突，并以赏赐、婚姻和仪式维持可见的等级。王权的力量不只是能否下令，更在于这些地方精英是否仍把王室的命令视为值得兑现的名分。

\begin{causalchain}[从东方据点到关系网络]
\term{结构性原因}是周初中心人力有限、东方政治关系复杂，不能只靠从关中远征来治理。\term{使能条件}是宗族、军功、礼器与祭祀能把地方精英编进共同语言。成周提供了较近的会集点，封国则提供日常执行者；二者结合，才形成\eventref{WZ-EAST-CHENGZHOU}{西周·营建成周}所代表的东方网络。\term{反馈回路}也在这里产生：王室每次对功劳的承认，都增强了受命者的资源和声望；资源越集中于地方，王室日后越要依赖协商，而不能只凭命令。
\end{causalchain}

\subsection{从册命到地方能力}

现存青铜器铭文经常记载王命、赏赐、服事和祭祀，说明名分并非抽象口号：它要通过反复的仪式和器物被保存、展示并传给后人。不过，这类材料大多由获得命令或赏赐的贵族制作，保存的是成功进入王室关系网的人如何叙述自己，较少留下未获承认者、普通服役者或地方对抗者的声音。

因此，东方网络不能只用“中央强、地方弱”或反过来的二分法理解。在\eventref{WZ-CHENGKANG-CONSOLIDATION}{西周·成康巩固}所代表的阶段，王室权威能够为封国提供共同的秩序语言，封国也把本地人口、土地和武力接入更广的体系。长远看，同一机制又让地方家族积累连续的祭祀传统、军事组织和资源调动能力。西周末年的危机并非由周初封国单独造成，却会在这些已具实力的地方网络中寻找盟友、兵员和退路。

\begin{debate}[“封建”能解释多少]
“封建”是后世整理周代制度时形成的概念，适合描述授土、建国和宗族政治彼此相连的现象，却不应当替代对具体地区的考察。研究者会结合金文中的册命和赏赐、传世文献的追述、聚落与墓葬的区域差异，讨论王室控制如何随地点和时期改变。没有一件材料能单独列出完整封国清单，更不能据此断言所有封国都以同一种方式服从王室。
\end{debate}
